\chapter{奇数下半带：一组共同反射}\label{ch:odd}

\section{本讲的构造任务}
现在进入临界带，设 $n=2p+r$、$1\le r<p$ 且 $r$ 为奇数。
入口已由第\ref{sec:compression}讲解决，收尾已由第\ref{sec:blocks}讲解决。
因此只需在核心中取得第\ref{sec:controllers}讲的共同终端条件。

本讲完整展示一次典型的构造过程：先设计两组三步反射，
令它们的商成为三角矩阵，再同换块作用作交换子，产生一条完整剪切。
另一个两步反射提供横向方向。整个构造覆盖所有奇数余量，包括最紧的 $r=1$。

为了使交换两个等大块变得简单，取对称坐标
\[
u=z,\qquad v=a-\frac r2z,
\qquad A=\frac r2u+v,\quad B=\frac r2u-v,\quad C=-pu.
\]
交换 $A,B$ 的名称就是 $S=\diag(1,-1)$。
分母 $2$ 目前只是分析坐标的选择；构造本身仍由整数个位置组成。

\begin{proposition}[奇数余量的全部合法核心]\label{prop:odd}
若 $p\ge3$ 为奇素数、$1\le r<p$ 为奇数，则全部合法 $K_{p,r}$ 可归零。
\end{proposition}
\begin{proof}

\par\noindent\textbf{第一步：选择两组库存内的反射。}
采用 $u=z,v=a-rz/2$。写 $p-r=4k+2\epsilon$，$\epsilon\in\{0,1\}$。取 $\alpha_0=-r,\alpha_1=4-r$。第一步从 $A,B$ 各取 $(p+\alpha)/2,(p-\alpha)/2$ 份，产生 $p$ 份 $D$，保留 $r$ 份作载体；第二步按
\[
(x,y,c,v_D)=\left(\frac{p+2r-\alpha-2\epsilon}{4},\frac{p+\alpha-2\epsilon}{4},\epsilon,\frac{p-r}{2}\right)
\]
选择输出，第三步平均其 $p$ 元补集。两组输出分别为 $(r+k,k,\epsilon,2k+\epsilon)$ 和 $(r+k-1,k+1,\epsilon,2k+\epsilon)$；剩余库存分别为 $(r+2k+\epsilon,2k+\epsilon,r,4k+2\epsilon)$ 和 $(r+2k+\epsilon-2,2k+\epsilon+2,r,4k+2\epsilon)$。唯一紧的条件是 $k+\epsilon\ge1$，由 $p-r>0$ 保证。

这三次操作的顺序是：先产生一个新 $p$ 块，留出其中的 $r$ 份作载体，
再平均一个输出组，最后平均其补组。补组自动有 $p$ 个位置，
因此只需逐项检查所列输出组不超过剩余库存。
返回的实际尺度为 $p^{-2}$，矩阵为
\[
J_\alpha=\begin{pmatrix}-rp/2&-\alpha\\pn(r-2\epsilon)/4&rp/2\end{pmatrix},
\quad J_\alpha^2=pF_\alpha I,\quad F_\alpha=\frac{pr^2-n(r-2\epsilon)\alpha}{4}.
\]

\par\noindent\textbf{第二步：标量周期与参数环。}
平方标量模 $n$ 为 $p^4$，故非零且安全。令 $c_\epsilon=r$（$\epsilon=0$）或 $1$（$\epsilon=1$），$f_i=F_{\alpha_i}/c_\epsilon$。若 $\epsilon=0$，则 $f_0=r(3p+r)/4,f_1=f_0-n$；若 $\epsilon=1$，则 $f_0=r[pr+n(r-2)]/4,f_1=f_0-n(r-2)$。这给出
\[
\gcd(f_0,f_1)=\gcd(f_0f_1,n)=1,\qquad r\mid f_0,
\]
且二者奇偶相反。第二种情形的互素性由模 $n$、模 $r-2$ 的等式 $4f_0\equiv pr^2$ 得出。

相邻标量互素，故三角商的倍率 $f_1/f_0$ 约分后仍保留两者的全部素因子。
而 $f_0,f_1$ 奇偶相反、$r\mid f_0$，所以 $2$ 与 $r$ 都会成为参数环的单位。
在 $R=\Z[1/|f_0f_1|]$ 中，$2,r$ 为单位，故 $R$ 的单位无限。令 $S=\diag(1,-1)$，则
\begin{gather*}
T=J_{\alpha_1}J_{\alpha_0}^{-1}
=\begin{pmatrix}f_1/f_0&-2r/(c_\epsilon f_0)\\0&1\end{pmatrix},\\
[S,T]=U\left(\frac{4r}{c_\epsilon f_0}\right).
\end{gather*}
根参数为 $R$ 单位，引理\ref{lem:localize}给 $U(R)$。

\par\noindent\textbf{第三步：取得横向剪切。}
保留 $r$ 份旧 $B$，平均两组
\[
(A^{(p+r)/2},B^{(p-r)/2}),\qquad
(A^{(p-r)/2},B^{(p-r)/2},C^r),
\]
得到
\[
K=\begin{pmatrix}-r/2&1\\rn/4&r/2\end{pmatrix},\quad K^2=r(p+r)I/2.
\]
若 $H=KU(2/r)$，则 $[S,H]=L(n)$。置 $P=\left(\begin{smallmatrix}1&-2\\0&n\end{smallmatrix}\right)$，有
\[
P^{-1}U(t)P=U(nt),\quad
P^{-1}KU(8(p+r)t/r)K^{-1}P=L(nt).
\]

这里换基矩阵的行列式为 $n$，必须计入一层分母损失。
引理\ref{lem:deep}先给 $\Gamma(n^3,R)$。对 $t\in R$ 选整数 $k\equiv t\pmod{n^2R}$，分解 $L(nt)=L(nk)L(n(t-k))$，得到 $L(nR)$，再得 $\Gamma(n^2,R)$。

\par\noindent\textbf{第四步：接入共同终端。}
换到 $(z,y)$ 的矩阵 $C=\left(\begin{smallmatrix}1&0\\n/2&1\end{smallmatrix}\right)$。选 $4\mid k$ 且 $k\equiv-p^{-2}\pmod n$，则 $CU(k)C^{-1}$ 整数且满足定理\ref{thm:terminal}的方向条件。至少一个整数反射的行列式为偶数，而条目公因子为奇数（其左下项 $-\alpha_i$ 奇）；引理\ref{lem:content}提供有限标量 $2I$。定理\ref{thm:terminal}完成证明。
\end{proof}

\section{一个可以手算的实例：五平均的十一位置}
取 $p=5,r=1$，则 $n=11$、$k=1,\epsilon=0$。
两组来源的 $\alpha$ 为 $-1,3$，对应
\[
J_{-1}=\begin{pmatrix}-5/2&1\\55/4&5/2\end{pmatrix},
\qquad
J_3=\begin{pmatrix}-5/2&-3\\55/4&5/2\end{pmatrix}.
\]
两者的平方分别是 $20I,-35I$，即 $f_0=4,f_1=-7$。
对第一组，先从两个五块各取 $2,3$ 份；留出一份新值后，
剩余四类库存为 $(3,2,1,4)$，输出取 $(2,1,0,2)$。
对第二组，来源改取 $4,1$ 份，库存为 $(1,4,1,4)$，输出取 $(1,2,0,2)$。
两次所选输出均为五项，且逐项不超过库存。

三角商的倍率为 $-7/4$，所以参数环为 $\Z[1/14]$。
这时根参数 $4r/(c_\epsilon f_0)=1$，完整的 $U(R)$ 随即得到。
这个例子说明：大定理中的有限局部化环，是从实际分组的标量平方里读出来的。
它并不是事先任意指定的有理系数域。


